\documentclass{article}

\usepackage{amssymb}
\usepackage{amsmath}
\usepackage{changepage}
\usepackage{titlesec}
\usepackage{mathtools}
\usepackage{graphicx}
\usepackage{pifont}
\usepackage{multirow}
\usepackage{amsthm}

\newcommand{\cB}{\mathcal{B}}
\newcommand{\cF}{\mathcal{F}}
\newcommand{\cH}{\mathcal{H}}
\newcommand{\cE}{\mathcal{E}}
\newcommand{\cS}{\mathcal{S}}
\newcommand{\E}{\mathbb{E}}
\newcommand{\N}{\mathbb{N}}
\newcommand{\Z}{\mathbb{Z}}
\newcommand{\Q}{\mathbb{Q}}
\newcommand{\R}{\mathbb{R}}
\newcommand{\Prb}{\mathbb{P}}
\newcommand{\Rp}{\R^{+}}
\newcommand{\Rb}{\overline{\R}}
\newcommand{\Rm}{\R^{-}}
\newcommand{\subelement}[1]{\setminus\{#1\}}
\newcommand{\subzero}{\subelement{0}}
\newcommand{\quantifyspace}{\hspace{2px} \big | \hspace{5px}}
\newcommand{\cheekynewline}{$\text{ }$\\}
\newcommand{\Tau}{\mathcal{T}}
\newcommand{\dstar}{\ding{65}}
\newcommand{\Pow}{\mathcal{P}}
\newcommand{\vp}{\varphi}
\newcommand{\oR}{\overline{\R}}
\newcommand{\wt}{\widetilde}
\newcommand{\M}{\mathcal{M}}
\newcommand{\df}[1]{\hspace{2 px} d #1 \hspace{2 px}}
\newcommand{\gral}[3]{\int_{#1} #2 \hspace{2 px} d #3 \hspace{2px}}
\newcommand{\lo}{\Bigg (}
\newcommand{\lc}{\Bigg )}
\newcommand{\dbl}{\Bigg | \Bigg |}
\newcommand{\PS}{(\Omega, \cF, \Prb)}
\newcommand{\FPS}{(\Omega, \cF, \{ \cF_n\}_{n \geq 0}, \Prb)}
\newcommand{\cFPS}{(\Omega, \cF, \{ \cF_t\}_{t \geq 0}, \Prb)}
\newcommand{\Fi}{\{\cF_n\}_{n \geq 0}}
\newcommand{\cFi}{\{\cF_t\}_{t \geq 0}}
\newcommand{\spc}{\hspace{7 px}}
\newcommand{\sspc}{\hspace{3 px}}
\newcommand{\vspc}{\vspace{7px}}
\newcommand{\la}{\langle}
\newcommand{\ra}{\rangle}
\newcommand{\ind}[1]{1_{\{#1\}}}
\newcommand{\norm}[2]{||#1||_{#2}}
\newcommand{\Var}{\text{Var}}
\newcommand{\ale}{\text{a.e.}}
\newcommand{\als}{\text{a.s.}}
\newcommand{\indep}{\perp\!\!\!\!\perp} 
\newcommand{\spdot}{\text{ } \cdot \text{ }}
\newcommand{\esssup}{\text{ess sup}}
\theoremstyle{definition}
\newtheorem{defn}{Definition}
\newtheorem{thm}{Theorem}
\newtheorem{lemma}{Lemma}
\newtheorem{cor}{Corollary}
\newtheorem{claim}{Claim}
\newtheorem{obj}{Objective}
\newtheorem{assume}{Assumption}

\graphicspath{ {./media/} }

\title{Optimal Stopping of Ornstein-Uhlenbeck Processes}

\author{Max David Feldman, B.S. \\ Supervisor:  Dr. Mihai Sîrbu}

\date{The University of Texas at Austin, 2026}

\setlength{\parindent}{0pt}

\begin{document}

\maketitle

\titleformat*{\section}{\bfseries}

% Abstract 
In this expository thesis we will start Part I by introducing some foundational concepts of measure-theoretic probability and the theory of continuous time martingales.
We will then make use of some analysis tools to define a concept essential to the general theory of optimal stopping, namely the essential supremum.
Part I will then finish with motivation and definition of the markov property.
In Part II we formulate and attack the general problem of optimal stopping and find its solution in terms of the Snell Envelope.
We finish Part II with motivation for what we hope the markov property will give us.
In Part III we look at a specific problem, namely that of valuing the perpetual American put option, and get a sense of what stands to be gained
from the markov property and additional structure in an application specific question of optimal stopping.
In Part IV we finally use the machinery and intuition we have built up to address the Ornstein-Uhlenbeck problem.

\section*{Part I: Preliminaries}
\subsection*{\quad I.I martingale theory}

\begin{defn} Given a continuous time filtered probability space $\cFPS$ and a $\cFi$-adapted stochastic process $(X_t)_{t \geq 0}$ 
We say that $X_t$ is a \textbf{submartingale} with respect to $\cFi$ if, \\
\[
    \forall t < s \spc \spc X_t \leq \E[ X_s \quantifyspace \cF_t] \spc  \als
\]
If the reverse inequality holds we say that $(X_t)_{t \geq 0}$ is a \textbf{supermartingale} with respect to $\cFi$. \\
If $(X_t)_{t \geq 0}$ is both a submartingale and a supermartingale with respect to $\cFi$ then we say that $(X_t){t \geq 0}$ is a \textbf{martingale} with respect to $\cFi$ \\
\end{defn}

\begin{defn}
    Two stochastic processes $(X_t)_{t \geq 0}$ and $(Y_t)_{t \geq 0}$ on $\PS$ are said to be \textbf{modifications} of each other if
    \[
        \Prb[X_t = Y_t] = 1 \spc \forall t \geq 0
    \]
    $X$ and $Y$ are further said to be \textbf{indistinguishable} if
    \[
        \Prb[X_t = Y_t \sspc \forall t \geq 0] = 1
    \]
\end{defn}

We now build some machinery that will be needed to prove the Upcrossing Inequality.

\begin{defn}
    Given $(X_t)_{t \geq 0}$ a stochastic process on $\PS$ and for a finite set of times $F \subset [s,t] \subset [0,\infty)$ and real valued constants $a < b$
    We define $T_1 := \min \{u \in F \quantifyspace X_u < a \}$ and $T_2 := \min \{u \in F \quantifyspace u > T_1, X_u > b\}$ \\
    Now we define inductively $T_{2k - 1} := \min \{u \in F \quantifyspace u > T_{2k - 2}, X_u < a\}$ and $T_{2k} := \min \{u \in F \quantifyspace  u > T_{2k-1}, X_u > b\}$ \\
    We then define $U_F (X, [a,b]) := \max \{k \quantifyspace T_{2k} \leq t\}$ \\
    And for any infinite set of times $I$ we define,
    \[
        U_I(X, [a,b]) := \sup_{F \subset I, F \text{ finite}} U_F(X, [a,b])
    \]
\end{defn}

\begin{lemma} Upcrossing inequality \\
    Given $(X_t)_{t \geq 0}$ a RC-submartingale and times $s < t$ we have,
    \[
        \E[U_{[s,t]}(X, [a,b])] \leq \Big ( \E[X_t^+] + |a| \Big ) (b - a)^{-1}
    \]
    This implies that $\Prb [U_{[s,t]}(X, [a,b]) = \infty] = 0$ \\

    \proof \cheekynewline
        We first prove that the result holds for $(M_n)_{n \geq 0}$ a discrete time sub-martingale over a finite time steps \\
        It will suffice to show that for fixed $N \in \N$ we have $\E [U_{\{0,1,..,N\}}(X, [a,b])] \leq \Big ( \E[M_N^+] + |a| \Big ) (b - a)^{-1}$ \\
        Now take $H_n := \sum_{j=0}^\infty \ind{T_{2j + 1} \leq n \leq T_{2j + 2}}$ \\
        And take $N_n := a \vee M_n = a + (M_n - a)^+$ since $a \vee \cdot$ is convex and non-decreasing we can apply conditional Jensen's and see that for $n \leq m$,
        \[
            N_n = a \vee M_n \leq a \vee \E[M_m \quantifyspace \cF_n] \leq \E[N_m \quantifyspace \cF_m]
        \]
        And so $N_n$ is a submartingale and since $1 - H_n \geq 0$ we have that $((1 - H) \cdot N)_n$ is a submartingale \\
        Also by construction, $(H \cdot N)_n \geq (b - a) U_{\{0, 1, .., n\}}([a,b])$
        % \textit{The above can be heuristically interpreted as a portfolio strategy of buying $N$ above $b$ and selling below $a$ is a submartingale, since for every passage from above $b$ to below $a$ we lose at least $b-a$ we can hope to bound the expected number of times this happens as each of them incurs loss} \\
        We find that for fixed $n \in \N$,
        \[
            0 = ((1 - H) \cdot N)_0 \leq \E[(1 \cdot N)_n - (H \cdot N)_n] 
        \]
        And so,
        \[
            (b - a) \E[U_{\{0,1..,n\}}(X, [a,b])] \leq \E[(H \cdot N)_n] \leq \E[N_n - N_0] =
        \]
        \[
            \E[(M_n - a)^+ - (M_0 - a)^+]
            \leq \E[(M_n - a)^+] \leq \E[M_n^+] + |a|
        \]
        This proves the discrete case. \\
        \\
        By choosing $\{q_n\}_{n = 1}^\infty$ to enumerate $\{b\} \cup (\Q \cap [s, t])$ we can take $F_n := \{q_1, .., q_n\}$, and by right continuity of $X_t$ we will get,
        \[
            U_{[s,t]}(X, [a,b]) = \lim_{n \to \infty} U_{F_n}(X, [a,b])
        \]
        Further restricting to finite times we find that $\E[U_{F_n}(X, [a,b])] \leq \E[X_t^+] + |a|$ from the discrete case, and since this holds for all $F_n$ by Fatou's lemma we find,
        \[
            \E \big [U_{[s,t]}(X, [a,b]) \big ] = \E \Big [ \liminf_{n \to \infty } U_{F_n}(X, [a,b]) \Big ] \leq \liminf_{n \to \infty} \E \big [ U_{F_n}(X, [a,b]) \big ] \leq \frac{\E[X_t^+] + |a|}{(b-a)}
        \]
        QED
        \\
        
\end{lemma}

% (MAYBE) include submartingale convergence  \\

\begin{defn}
    We are given a filtered probability space $\cFPS$. 
    An $\cF$-measurable random variable $T: \Omega \to [0,\infty]$ is called a random time. \\
    A random time $T$ is said to be a \textbf{stopping time} if $\{T \leq t\} \in \cF_t$ for all $t \geq 0$ 
\end{defn}

% (MAYBE) DEF Uniform Integrability \\

\begin{defn}
    A collection of random variables $(V_\alpha)_{\alpha \in A}$ we say that $(V_\alpha)$ is
    \textbf{Uniformly Integrable (U.I.)} if,
    \[
        \lim_{K \to \infty} \sup_{\alpha \in A} \sspc \E [V_\alpha \ind{|V_\alpha| \geq K}] = 0
    \]
\end{defn}

% THM Optional Sampling \\
\begin{thm}
    \textbf{Optional Stopping} \\
    Let us say that $(M_t)_{t \geq 0}$ and is a right continuous submartingale on $\cFPS$ and we are given optional times $S, T$ s.t. $0 \leq S \leq T$. \\
    If $(M_{t \wedge T}^+)_{t \geq 0}$ is U.I. then $M_S \leq \E[M_T \quantifyspace \cF_{S^+}]$
\end{thm}

\begin{thm}
    A (super) submartingale $(X_t)_{t \geq 0}$ with the property that the map $t \mapsto \E[X_t]$ is right continuous has an RCLL modification. \\
\end{thm}

\begin{thm}
    \textbf{Doob's Maximal Inequality} \\
    Given $(M_t)_{t \geq 0}$ a submartingale we define $M_t^* := \sup_{0 \leq s \leq t} M_t^+$ and then for $p \in (1,\infty)$ we have,
    \[
        \E[(M_t^*)^p] \leq \Big ( \frac{p}{p-1} \Big )^p  \sspc \E[(M_t^+)^p]
    \]
\end{thm}

\begin{defn}
    We say that a filtration $\cFi$ for a probability space $\PS$ satisfies the \textbf{usual conditions} if $\mathcal{W} := \{A \in \cF \quantifyspace \Prb[A] = 0\}$ is contained in $\cF_0$
    and $(\cF_t)_{t \geq 0}$ is right continuous, i.e.
    \[
        \forall t \geq 0 \spc \spc \cF_t = \cF_t^+ = \bigcap_{s > t} \cF_s
    \]
\end{defn}

\begin{defn}
    Given $(M_t)_{t \geq 0}$ RC and adapted on $\cFPS$ we say that $(M_t)_{t \geq 0}$ is a \textbf{local martingale} if there exist a sequence of stopping times $\{T_n\}_{n =1}^\infty$ s.t. $T_n \nearrow \infty$ a.s. and for all $n \in \N$ we have that $(M_{t \wedge T_n})_{t \geq 0}$ is a martingale. \\
\end{defn}

\begin{defn} (KS)
    A stochastic process $(X_t)_{t \geq 0}$ is said to be a \textbf{semimartingale} if there exist a local martingale $(M_t)_{t \geq 0}$ starting at 0 and a finite variation adapted process $(A_t)_{t \geq 0}$ starting at 0 s.t.
    \[
        X_t = X_0 + M_t + A_t
    \]
\end{defn}

\begin{thm}
    Let $M,N$ be continuous path local martingales then there exists a unique process $\la M, N \ra_t$ that starts at 0 and s.t. $(M_t N_t - \la M, N \ra_t)_{t \geq 0}$ is a local martingale.
    Also $\la M, N \ra$ is continuous and has bounded variation paths.
    We refer to $\la M, N \ra$ as the \textbf{quadratic variation} of $M$ against $N$
\end{thm}

\begin{cor}
    We note that for a local martingale $M$ we must have that $\la M, M \ra$ is non-decreasing
\end{cor}

\begin{defn}
    A stochastic process $(B_t)_{t \geq 0}$ starting at 0 and with continuous paths on an event of probability 1 is said to be a \textbf{1-dimensional brownian motion (brownian motion)} if
    \begin{enumerate}
        \item{(i)} for all $s < t$ we have that $(B_t - B_s)$ has a normal distribution with mean $0$ and $\sigma^2 = t-s$ \\
        \item{(ii)} for all finite collections of deterministic times $0 \leq t_0 < t_1 < t_2 < t_3 < .. < t_n$ we have that the collection of increments as random variables $\{B_{t_k} - B_{t_{k-1}}\}_{k=1}^n$ is independent
    \end{enumerate}
\end{defn}

\begin{lemma}
    There exists a probability space $\PS$ s.t. there exists a brownian motion on this space.
\end{lemma}

\begin{lemma}
    \textbf{Itô's Lemma} \\
    Given $(X_t)_{t \geq 0}$ a continuous semi martingale and $f \in C^2(\R)$ we have,
    \[
        \Prb \bigg [ f(X_t) = f(X_0) + \int_0^t f'(X_u) \df{X_u} + \frac{1}{2} \int_0^t f''(X_u) \df{\la X, X \ra_u} \spc \forall t \in \R_+ \bigg ] = 1
    \]
\end{lemma}
\cheekynewline
Infact Itô's lemma extends to a broader class of $f$.

\begin{thm}
    (KS Ch 3.7 problem 7.3)
    Given $(X_t)_{t \geq 0}$ a continuous semi martingale and $f \in C^1(\R)$ with $f'$ being absolutely continuous then $f''$ exists $\lambda$-a.e.
    and we have,
    \[
        f(X_t) = f(X_0) + \int_0^t f'(X_u) \df{X_u} + \frac{1}{2} \int_0^t f''(X_u) \df{\la X, X \ra_u} \spc
        \forall t \in \R_+ \spc \spc \Prb - \als
    \]
\end{thm}

\begin{thm}
    \textbf{Paul Levy Characterization of Brownian Motion} (A.S. Blog) \\
    Given a local martingale $(X_t)_{t \geq 0}$ starting at 0 on the filtered probability space $\cFPS$ the following are equivalent:
    \begin{enumerate}
        \item[(i)] $X_t$ is a Brownian motion on specified filtered probability space
        \item[(ii)] $X$ is continuous and $X_t^2 - t$ is a local martingale
        \item[(iii)] $\la X, X \ra_t = t$ for $t \in \R_+$ 
    \end{enumerate}
\end{thm}

\subsection*{\quad I.II essential supremum}

\begin{defn}
Consider a probability space $\PS$ and $\mathcal{Q}$ a nonempty collection of non-negative random variables defined on $\PS$.
The \textbf{essential supremum} of $\mathcal{Q}$ (denoted ess sup $\mathcal{Q}$) is a random variable $X^*$ with the following two properties 

\begin{enumerate}
    \item $\forall X \in \mathcal{Q}, \spc X \leq X^* \spc \als$
    \item $(X \leq Y \sspc \als \sspc \forall X \in \mathcal{Q}) \Rightarrow (X^* \leq Y \sspc \als)$
\end{enumerate}
\end{defn}

\begin{defn}
    Let $A \in \cF$ and say $k \in \N$ where $A_1, A_2, .., A_k \in \cF$ comprise a partition of $A$. And say $X_1, X_2, .., X_k \in \mathcal{Q}$
    We then say that $\pi = (k, A_1, A_2, .., A_k, X_1, X_2, .., X_k)$ is a \textbf{$\mathcal{Q}$-partition of $A$}. \\
    We then let $P_A$ denote the set of all $\mathcal{Q}$-partitions of $A$. \\
\end{defn}

Now we construct $\mu_\pi^\infty: \cF \to \R_+$ as follows,
\[
    A \xmapsto{\mu_\pi^\infty} \E \bigg [ \sum_{j=1}^k X_j 1_{A_j} \bigg ]
\]
We then define $\mu^\infty: \cF \to \R_+$ as,
\[
    A \xmapsto{\mu^\infty} \sup_{\pi \in P_A} \mu_\pi^\infty(A)
\]

\begin{lemma}
    $\mu^\infty: \cF \to \R_+$ is a measure on $\Omega$ and $\mu^\infty$ is absolutely continuous with respect to $\Prb$.
\end{lemma}

\vspc

\begin{thm}
    Consider $\mathcal{Q}$ a collection of random variables defined on the probability space $\PS$.
    Then ess sup $\mathcal{Q}$ exists and is unique up to $\Prb$-a.s. equivalence.

    \proof \cheekynewline
    \\
    First we show that if the essential supremum of $\mathcal{Q}$ exists then it is unique up to $\Prb$-a.s. equivalence \\
    This follows trivially from the second property of the essential supremum. \\
    Say that $A$ and $B$ both satisfy the definition of essential supremum of $\mathcal{Q}$ \\
    Since $X \leq A$ a.s. and $X \leq B$ a.s. for all $X \in \mathcal{Q}$ we then have that $B \leq A$ a.s. and $A \leq B$ a.s. thus $A = B$ a.s. \\
    \\
    Now we show existence ... \\
    Recall that $\mu^\infty << \Prb$ and now define $X^* := \frac{\df{\mu^\infty}}{\df{\Prb}}$ i.e. $X^*$ is the Radon-Nikodym derivative of $\mu^\infty$ w.r.t. $\Prb$ \\
    We claim that $X^*$ satisfies the criterion of essential supremum of $\mathcal{Q}$ \\
    \\
    Fix some $X \in \mathcal{Q}$ \\
    Now fix a test event $A \in \cF$ and consider the $\mathcal{Q}$-partition of $A$, $\pi^* := (1, A, X)$ and observe,
    \[
        \E \big [X 1_A \big ] = \mu_{\pi^*}^\infty(A) \leq \sup_{\pi} \mu_\pi^\infty(A) = \mu^\infty(A) = \E \Big [\frac{\df{\mu^\infty}}{\df{\Prb}} 1_A \Big ] = \E \big [X^* 1_A \big ]
    \]
    Since this holds for all $A \in \cF$ we have that $X \leq X^*$ a.s and we have shown the first property. \\
    \\
    Now we fix some random variable $Y$ s.t. $X \leq Y$ a.s. for all $X \in \mathcal{Q}$ \\
    Observe,
    \[
        \E \big [X^* A \big ] =
        \sup_{\pi = (k, A_1 .. A_k , X_1 .. X_k) \in P_A } \mu_\pi^\infty (A) =
        \sup_{\pi \in P_A} \E \bigg [\sum_{j=1}^k X_j 1_{A_j} \bigg ]  \leq
        \sup_{\pi \in P_A} \E \bigg [ Y \sum_{j=1}^k 1_{A_j} \bigg ] = \E \big [ Y 1_{A} \big ] 
    \]
    This shows the second condition. \\
    QED
\end{thm}

\begin{lemma}
    Consider $\mathcal{Q}$ a collection of random variables defined on the probability space $\PS$ s.t. $\mathcal{Q}$ is closed under pairwise maximum i.e.
    $A,B \in \mathcal{Q}$ implies $A \vee B \in \mathcal{Q}$, then there is a non-decreasing sequence $\{C_n\}_{n=1}^\infty \subset \mathcal{Q}$ s.t. $X^* = \lim_{n \to \infty} C_n$ a.s.
\end{lemma}

\subsection*{ \quad I.III markov property and strong markov property \& Ornstein-Uhlenbeck processes}
\cheekynewline
We introduce the idea of the Markov Property which is motivated by the idea that for some processes the increments are dependent only on state at beginning of the start of the time interval of the increment and not on the entire trajectory. \\
\begin{defn}
    Given a state space $(E, \mathcal{E})$ a \textbf{transition function} is a function \\
    $P: [0,\infty) \times E \times [0,\infty) \times \mathcal{E} \to [0,1]$ s.t.
    For each $s,t \in [0,\infty), e \in E$ we have that $P(s, e; t, \cdot)$ is a probability measure and for every $0 \leq s < r < t$ and $e \in E$ we have,
    \begin{equation}
        P(s, e; t, A) = \int_E P(s,e;r, dy) P(r, y; t, A)
        \tag{\textsc{Chapman-Kolmogorv equation}}
    \end{equation}
    If in addition for all $s,t \geq 0$ and $\alpha > 0$ we have \\
    $$P(s, \spdot; s+\alpha, \spdot) = P(t, \spdot; t + \alpha, \spdot)$$ \\
    then we say that $P$ is \textbf{time-homogenous} and define \\
    $$P(\alpha, \spdot; \spdot) := P(s, \spdot; s+\alpha, \spdot) = P(t, \spdot; t+\alpha, \spdot)$$
\end{defn}

\begin{defn}
    Given $(X_t)_{t \geq 0}$ as $(E, \mathcal{E})$-valued random elements on the filtered probability space $\cFPS$ s.t. $X_t \in \cF_t$ for all $t \geq 0$
    and given a transition function $P$
    we say that $X_t$ is a \textbf{Markov Process} with respect to $\cFi$ with transition function $P$ if for all $s < t$ and measurable functions $f: E \to \R_+$ we have 
    \begin{equation*}
        \E[f(X_t) \quantifyspace \cF_s] = \int_E P(s, X_s; t, dy) f(y) \spc \spc \Prb-\als
    \end{equation*}
    If in addition $P$ is time-homogeneous then we say that $X_t$ is a \textbf{time-homogeneous} Markov Process. \\
\end{defn}
\cheekynewline

\begin{defn}
    Given $(X_t)_{t \geq 0}$ a $(E, \mathcal{E})$-valued time-homogeneous Markov process with respect to $\cFi$ with transition function $P$
    we say that $X_t$ satisfies the \textbf{Strong Markov Property} if for all measurable functions $f: E \to \R_+$ and all stopping times $\tau$ we have,
    \[
        \E[f(X_{\tau+\alpha}) \quantifyspace \cF_\tau] = \int_E P(\alpha, X_\tau; dy) f(y) \spc \spc \Prb-\als \text{ on } \{\tau < \infty\}
    \]
    \small{ \spc we note that some authors relax away the requirement that $X$ must be time-homogenous in order to have the strong markov property}
\end{defn}


We now introduce a specific instance of a class of Markov Processes,
\begin{defn}
    Given $(X_t)_{t \geq 0}$ an adapted stochastic process on the filtered probability space $\cFPS$
    we say that $X_t$ is an \textbf{Ornstein-Uhlenbeck process} if $X_t$ satisfies the following SDE.
    \[
        \df{x_t} = \theta (\mu - x_t) \df{t} + \sigma \df{W_t}
    \]
    Where $W_t$ is a Wiener process and $\theta, \sigma > 0$ and $\mu \in \R$ are constants \\
\end{defn}

\section*{Part II: Optimal Stopping}
\subsection*{\quad II.I Formulation of general problem}
We are given a filtered probability space $\cFPS$ s.t. the filtration satisfies the usual conditions and we specifically impose $\cF_0 = \mathcal{W}$ \\
We are also given a time horizon $T \in (0, \infty]$ \\
We are given $(Y_t)_{t \geq 0}$ a non-negative valued and right continuous stochastic process on our probability space. \\
In the case that $T = \infty$ we consider $\cF_\infty := \sigma ( \cup_{t \geq 0} \cF_t)$ and $Y_\infty := \limsup_{t \to \infty} Y_t$ \\
We now define $\mathcal{S}$ to be the set of all $\cFi$ stopping times. \\
Now we take
\[
    \cS_\tau := \{\theta \in \cS \quantifyspace \theta \geq \tau \sspc \als \}
\]
\cheekynewline

It is useful for intuition to conceive of $Y_t$ as representing the value that exercise of a particular derivatives contract (for now just a single exercise derivative) at time $t$ yields.

\begin{obj}
    We want to achieve the following
    \begin{enumerate}
        \item[(i)] Compute the maximal expected value of stopping $Y$
            \[
                Z_0 := \sup_{\tau \in \cS} \E[Y_\tau]
            \]
        \item[(ii)] Find conditions that yield existence of a stopping time $\tau_*$ that achieves the value in the prior sub-objective
        \item[(iii)] Build understanding of $\tau_*$
    \end{enumerate}
\end{obj}

\subsection*{\quad II.II Snell Envelope}
The essential (pun intended) piece of intuition we need to start to build the right machinery
is that given two stopping times $\tau, \theta \in \cS$ we can construct another stopping time $\gamma$ which takes values of $\tau$ when preferable and values of $\theta$ when preferable. \\
Specifically, we define $A := \{\E[Y_\tau - Y_\theta \quantifyspace \cF_{\tau \wedge \theta}] \geq 0\}$ and then define $\gamma := \tau 1_A + \theta 1_{A^c}$ \\
This observation tell us that
given a fixed $\tau \in \cS$ we have that $\{\E[Y_\theta \quantifyspace \cF_\tau] \}_{\theta \in \cS_\tau}$ is closed under pairwise maximization. \\
At a glance one might hope that this family would have a maximal element but we must recall that each of these elements are only defined up to a.s. equality so if we tried to make sense of such an $\omega$ by $\omega$ supremum we could get the random variable $\omega \mapsto \infty$. \\
This should be enough to tell you that the essential supremum is the right tool, we can now get to the main definition \\
\\
\begin{defn}
    We define the following,
    \[
        Z_v := \esssup_{\tau \in \cS_v} \E[Y_\tau \quantifyspace \cF_v]
    \]
    A modification of this will be the "Snell Envelope"
\end{defn}

We only consider the non-trivial case where
\[
    0 < Z_0 = \sup_{\tau \in \cS} \E[Y_\tau] < \infty
\]
and so we will assume this going forward \\
\cheekynewline

\begin{assume}
    We now impose the condition that
    \begin{equation}
        \E[\sup_{t \in [0,T]} Y_t] < \infty
    \end{equation}
    using this condition we can achieve sub-objective (ii).
\end{assume}

\begin{lemma}
    For all $v \in \cS$ and $\tau \in \cS_v$ the family $\{\E[Y_\rho \quantifyspace \cF_v]\}_{\rho \in \cS_v}$ is closed
    under pairwise maximization. Furthermore, there is a sequence $\{\rho_n\}_{n=1}^\infty$ of stopping times in $\cS_\tau$ s.t.
    the sequence $\{\E [Y_{\rho_n} \quantifyspace \cF_v] \}_{n=1}^\infty$ is non-decreasing and
    \[
        \esssup_{\rho \in \cS_\tau} \E[Y_\rho \quantifyspace \cF_v] = \lim_{n\to \infty} \E[Y_{\rho_n} \quantifyspace \cF_v]
    \]
    \proof
    The first statement of this lemma is exactly our observation at the beginning of this Snell Envelope section and the second statement is a direct result of the first statement and Lemma 5 of the Essential Supremum section. \\
\end{lemma}

The reader may note that $Z_v$ is a random variable that has been defined as a function of a stopping time $v \in \cS$, it is not clear that this definition of $Z$ is consistient i.e. if we consider the process given by $Z$ at deterministic times we do not know that stopping this process at $v$ is the same as what we defined as $Z_v$, we address this issue now and build some understanding of the relation between $Z$ and $Y$.
\begin{claim} ***
    For all $v, \sigma \in \cS$ and $\tau \in \cS_v$ we have
    \[
        Z_v = Z_\sigma \spc \als \text{ on } \{\sigma = v\}
    \]
    \[
        \E[Z_\tau \quantifyspace \cF_v] = \esssup_{\rho \in \cS_\tau} \E[Y_\rho \quantifyspace \cF_v] \spc \als
    \]
    \[
        \E[Z_\tau \quantifyspace \cF_v] \leq Z_v \spc \als
    \]
    \[
        \E[Z_\tau] = \sup_{\rho \in \cS_\tau} \E[Y_\rho] \leq Z_0 < \infty
    \]

    \proof ...
\end{claim}

\begin{defn}
    For a given $v \in \cS$ we introduce $\cS_v^* := \{\tau \in \cS_v \quantifyspace \tau > v \sspc \als \text{ on } \{v < T\} \}$ and we define,
    \[
        Z^*_v := \esssup_{\tau \in \cS_v^*} \E[Y_\tau \quantifyspace \cF_v]
    \]
\end{defn}
Observe that our proof of Claim 1 also holds if we replace $Z$ with $Z^*$ so $Z^*$ has all the outlined properties. \\
We now want to show right continuity of $Z^*$,
\begin{claim}
    For any $v \in \cS$ and a sequence $v_n \subset \cS_v^*$ s.t. $\lim_{n \to \infty} v_n = v$ a.s. we have,
    \[
        Z_v = Z_v^* \vee Y_v \spc \als
    \]
    \[
        \int_A Z_v^* \df{\Prb} = \lim_{n \to \infty} \int_A Z_{v_n}^* \df{\Prb} \spc \spc \forall A \in \cF_v
    \]
\end{claim}

\begin{cor}
    For any $v \in \cS$, we have $Z_v^* = Z_v$ a.s.
\end{cor}

\begin{defn}
    Since $Z$ (as a function from stopping times to random variables) is consistent in the sense of Claim 1 we take the deterministic times and consider the supermartingale $(Z_t)_{t \geq 0}$.
    Since $t \mapsto \E [Z_t^*] = \E[Z_t]$ is right continuous we may apply Theorem 2 to see that there exists a supermartingale $(Z^0_t)_{t \geq 0}$ with RCLL paths that is a modification of $Z$.
    We call $Z^0$ the \textbf{Snell Envelope} of $Y$.
\end{defn}

\begin{thm}
    Given a stopping time $\tau \in \cS$ we define the following,
    \[
        \theta_\tau^* := \inf \{t \geq \tau \quantifyspace Z^0_t = Y_t \} \wedge T
    \]
    We claim that $\E[Y_{\theta_\tau^*}] = \sup_{\rho \in \cS_\tau} \E[Y_\rho] $ i.e. $\theta_\tau^*$ is the optimal stopping time beginning from time $\tau$.
\end{thm}

\begin{thm}
    Any stopping time $\tau_*$  is optimal starting from time 0 (meaning $\E[Y_{\tau_*}] = Z^0_0 = \sup_{\rho \in \cS} \E[Y_\rho]$) if and only if
    \[
        Z^0_{\tau_*} = Y_{\tau_*} \spc \als 
    \]
    and the stopped supermartingale
    \[
        \big ( Z^0(t \wedge \tau_*) \big )_{t \geq 0} \text{ is a martingale.}
    \]
\end{thm}

\begin{defn}
    We say that a stochastic process $(X_t)_{t \geq 0}$ is \textbf{class D} if it is U.I. \\
    And we say that $(X_t)_{t \geq 0}$ is \textbf{class DL} if for any fixed $t$ we have that $\{X_\tau\}_{\tau \leq t, \tau \in \cS}$ is U.I. \\
\end{defn}

\begin{thm}
    \textbf{Doob-Meyer Decomposition Theorem} (KS Theorem 1.4.10) \\
    Say we are given $(H_t)_{t \geq 0}$ a right continuous submartingale on a filtered probability space with filtration satisfying the usual conditions. \\
    Then if $(H_t)_{t \geq 0}$ is class DL we can decompose $H$ as follows,
    \[
        H_t = M_t + A_t \spc \spc \forall t \geq 0
    \]
    Where $(M_t)_{t \geq 0}$ is an RCLL martingale and $(A_t)_{t \geq 0}$ is an adapted RCLL non-decreasing process. \\
    If in addition $H$ is class D then $(M_t)_{t \geq 0}$ is U.I. and $(A_t)_{t \geq 0}$ has an integrable terminal element.
\end{thm}

\begin{cor}
    Using the condition we imposed in Assumption 1 we get that $Y$ and in turn $Z$ and $Z^0$ are in classes DL and D.
    We can now apply the Doob-Meyer Decomposition Theorem to get a decomposition $Z^0_t = M_t - A_t$ where $(M_t)_{t \geq 0}$ is a UI martingale and $(A_t)_{t \geq 0}$ is a non-decreasing adapted process with an integrable terminal element. \\
\end{cor}

\subsection*{\quad II.III The Markov Dream}
At a high level we can think of the (strong) Markov property for a process $(X_t)_{t \geq 0}$ as the property that increments starting from a time $(\tau)$ $t$ are dependent only on the current state and not on the entire trajectory up to time $t$. \\
We now consider $(X_t)_{t \geq 0}$ having the strong markov property and with state space $(\R, \cB(\R))$. 
Given our high level description of the (strong) Markov property it may seem intuitive that the optimal stopping of $(X_t)_{t \geq 0}$ would depend only on the state at time of stoppage i.e. that there must be a set of states $S \in \cB$ s.t. $\tau_* := \inf \{ t \geq 0 \quantifyspace X_t \in S \}$ is optimal \\

\section*{Part III: American Put Option}

\section*{Part IV: Stopping magnitude of Ornstein-Uhlenbeck with linear time penalty}
Example comes from Example 10.1 in Chapter 10.2 of Shiriaev and Peskir


\subsection*{Alternate: 11.2 Space Change example is simple}

\subsection*{Alternate: 12.2 Measure Change is between 11.2 and ex 10.1 in terms of detail}
. \\


\section*{References}


Doob Meyer Decomp Proof [https://arxiv.org/pdf/1012.5292]

\end{document}
